% =============================================================================
\section{Finite State Machine Design}
\label{sec:fsm}
% =============================================================================

\subsection{Formal Mathematical Specification}

To guarantee mathematical determinism, absence of deadlocks, and verifiable safety compliance, the
washing machine controller is formalized as an \textbf{Extended Finite State Machine (EFSM)}.
Formally, the control unit is defined by the 6-tuple:
\begin{equation}
  \mathcal{M} = \big\langle \mathcal{S}, \Sigma, \Lambda, s_0, \delta, \omega \big\rangle
  \label{eq:fsm-tuple}
\end{equation}
operating over an extended context space $\mathcal{C}$.

\subsubsection{Constituent Sets and Context Spaces}

\begin{enumerate}[leftmargin=*]
  \item \textbf{Qualitative State Set ($\mathcal{S}$):} A discrete set of six operational states:
        \begin{equation}
          \mathcal{S} = \{\text{\st{STANDBY}}, \text{\st{COLLECTING}}, \text{\st{READY}}, \text{\st{RUNNING}}, \text{\st{PAUSED}}, \text{\st{ERROR}}\}
        \end{equation}
  \item \textbf{Extended Context Space ($\mathcal{C}$):} A tuple of bounded numerical variables:
        \begin{equation}
          \mathcal{C} = \mathbb{N} \times [0, T_{\text{cycle}}] \times \{0, 1\} \times [0, T_{\text{double}}] \times \mathcal{F}
        \end{equation}
        where $c = \langle b, t, n, w, f \rangle \in \mathcal{C}$ encapsulates:
        \begin{itemize}
          \item $b \in \mathbb{N}$: Accumulated coin balance in cents ($0 \le b \le 9999$).
          \item $t \in [0, 1800]$: Remaining cycle time budget in seconds ($T_{\text{cycle}} = 1800\,\text{s}$).
          \item $n \in \{0, 1\}$: Registered \code{STOP} press counter within the active window.
          \item $w \in [0, 1500]$: Residual time of the double-press window in milliseconds ($T_{\text{double}} = 1500\,\text{ms}$).
          \item $f \in \mathcal{F} = \{0, \dots, 7\}$: Bitmask of asserted hardware fault codes.
        \end{itemize}
  \item \textbf{Input Event Alphabet ($\Sigma$):} A discrete alphabet of ten mutually exclusive events:
        \begin{equation}
        \begin{aligned}
          \Sigma = \big\{ & \text{COIN}_{10}, \text{COIN}_{20}, \text{COIN}_{50}, \text{RUN}, \text{PAUSE}, \\
                          & \text{STOP}, \text{TICK}_{1\text{s}}, \text{TICK}_{1\text{ms}}, \text{FAULT}, \text{CLEARED} \big\}
        \end{aligned}
        \end{equation}
  \item \textbf{Output Alphabet ($\Lambda$):} The Cartesian product of indicator, actuator, and notification domains:
        \begin{equation}
          \Lambda = \Lambda_{\text{RLED}} \times \Lambda_{\text{BLED}} \times \Lambda_{\text{motor}} \times \Lambda_{\text{pump}} \times \Lambda_{\text{door}} \times \Lambda_{\text{return}}
        \end{equation}
  \item \textbf{Initial Configuration:} $s_0 = \text{\st{STANDBY}}$ with context $c_0 = \langle 0, 0, 0, 0, 0 \rangle$.
  \item \textbf{Transition Function ($\delta$):} $\delta: (\mathcal{S} \times \mathcal{C}) \times \Sigma \to (\mathcal{S} \times \mathcal{C})$.
  \item \textbf{Output Function ($\omega$):} $\omega: (\mathcal{S} \times \mathcal{C}) \times \Sigma \to \Lambda$.
\end{enumerate}

\subsubsection{Machine Classification: Hybrid Moore-Mealy Model}

The machine $\mathcal{M}$ is architecturally classified as a \textbf{Hybrid Moore-Mealy Machine}:
\begin{itemize}
  \item \textbf{Moore Property (Steady-State Outputs):} The optical status LEDs (\code{RLED}, \code{BLED}),
        the door deadbolt latch state, and the primary continuous actuator commands (\code{Motor}, \code{Pump})
        are strictly invariant within a given qualitative state $s \in \mathcal{S}$ (and phase partitioned within $t$).
        They depend solely on the current state $s$, satisfying $\omega_{\text{Moore}}: \mathcal{S} \to \Lambda_{\text{steady}}$.
  \item \textbf{Mealy Property (Transient Transition Actions):} Side effects such as clearing the coin balance
        ($b := 0$), dispensing pre-run refunded coins, and arming the double-press window timer
        ($w := 1500\,\text{ms}$) are instantaneous transient actions executed strictly upon the edge of transition
        firing, satisfying $\omega_{\text{Mealy}}: (\mathcal{S} \times \mathcal{C}) \times \Sigma \to \Lambda_{\text{transient}}$.
\end{itemize}

\subsection{State Transition Diagrams}

Figure~\ref{fig:fsm-main} depicts the operational state transitions under nominal working conditions.
Figure~\ref{fig:fsm-error} details the fail-safe transitions entering and recovering from the \st{ERROR} state.
Transition identifiers T$xx$ correspond bijectively to Table~\ref{tab:transitions}.

\tikzset{
  state/.style={circle, draw, thick, minimum size=2.3cm, align=center, font=\small\bfseries},
  tr/.style={-{Stealth[length=2.6mm]}, thick},
  lbl/.style={font=\footnotesize, align=center, fill=white, inner sep=2.5pt},
}

\begin{figure}[H]
  \centering
  \resizebox{\linewidth}{!}{%
  \begin{tikzpicture}
    % Discrete FSM States (generously spaced: S at (1,0), C at (7.5, 3.6), R at (15.5, 3.6), U at (15.5, -3.2), P at (7.5, -3.2))
    \node[state, fill=stIdle]  (S) at (1.0, 0.0)   {STANDBY\\[-1pt]\footnotesize\mdseries R:ON B:OFF};
    \node[state, fill=stIdle]  (C) at (7.5, 3.6)   {COLLECTING\\[-1pt]\footnotesize\mdseries R:ON B:OFF};
    \node[state, fill=stReady] (R) at (15.5, 3.6)  {READY\\[-1pt]\footnotesize\mdseries R:OFF B:ON};
    \node[state, fill=stRun]   (U) at (15.5,-3.2)  {RUNNING\\[-1pt]\footnotesize\mdseries R:OFF B:1Hz};
    \node[state, fill=stPause] (P) at (7.5,-3.2)   {PAUSED\\[-1pt]\footnotesize\mdseries R:OFF B:ON};

    % Power-on initialization
    \draw[tr] (-1.8, 0.0) -- node[lbl, above=2pt, pos=0.4]{Power-On} (S.west);

    % S -> C (T01) sloped above curve in open northern quadrant; high-arch C -> S (T06)
    \draw[tr] (S.north) to[bend left=16] node[lbl, sloped, above=2pt, pos=0.45]{T01: COIN [$<50$¢]} (C.west);
    \draw[tr, dashed] (C.north west) to[out=145, in=60, looseness=1.4] node[lbl, sloped, above=2pt, pos=0.5]{T06: 2$\times$STOP / Refund} (S.north west);
    
    % C -> R (T04)
    \draw[tr] (C.east) -- node[lbl, above=2pt]{T04: COIN [$\ge 50$¢]} (R.west);
    
    % Central corridor: S -> R (T02) routed into eastern quadrant (x ~ 11.5)
    \draw[tr] (S.east) .. controls (5.0, 0.8) and (10.5, 0.8) .. node[lbl, sloped, above=2pt, pos=0.7]{T02: COIN [$\ge 50$¢]} (R.south west);
    
    % Central corridor: R -> S (T06) routed below axis into western quadrant (x ~ 6.0)
    \draw[tr, dashed] (R.south west) .. controls (10.5, -0.6) and (5.0, -0.6) .. node[lbl, sloped, below=2pt, pos=0.4]{T06: 2$\times$STOP / Refund} (S.east);
    
    % Self loops for C and R (top, out=60, in=100)
    \draw[tr] (C) to[out=60, in=100, looseness=3.6] node[lbl, above=2pt]{T03: COIN [$<50$]} (C);
    \draw[tr] (R) to[out=60, in=100, looseness=3.6] node[lbl, above=2pt]{T07: COIN (surplus)} (R);

    % Cycle execution (R -> U vertical drop)
    \draw[tr] (R.south) -- node[lbl, right=6pt, pos=0.5]{T08: RUN / Start} (U.north);
    
    % RUNNING <-> PAUSED
    \draw[tr] (U.north west) to[bend right=22] node[lbl, above=2pt, pos=0.5]{T09: PAUSE} (P.north east);
    \draw[tr] (P.east) -- node[lbl, above=2pt, pos=0.5]{T14: RUN / Resume} (U.west);
    
    % Self loop for U (right side)
    \draw[tr] (U) to[out=20, in=-40, looseness=3.6] node[lbl, right=6pt]{T10: 1\,s Tick [$t > 1$]} (U);

    % Self loop for P (directed straight down at y=-5.0)
    \draw[tr] (P.south) to[out=-60, in=-120, looseness=3.5] node[lbl, below=2pt, pos=0.5]{T15: 1\,s Tick [$t > 1$]} (P.south);

    % P -> S return transition
    \draw[tr] (P.west) to[bend left=12] node[lbl, above=2pt, pos=0.5]{T16: $t=1$ \;|\; T13: 2$\times$STOP} (S.south east);

    % Deep bottom return loop U -> S
    \draw[tr] (U.south) .. controls (13.0, -7.0) and (2.5, -7.0) .. node[lbl, below=2pt, pos=0.5]{T11: $t=1$ \;|\; T13: 2$\times$STOP} (S.south);
  \end{tikzpicture}%
  }
  \caption{Formal state transition diagram: Nominal operation ($\text{T05}$ and $\text{T12}$ denote internal window arming)}
  \label{fig:fsm-main}
\end{figure}

\begin{figure}[H]
  \centering
  \resizebox{\linewidth}{!}{%
  \begin{tikzpicture}
    % Discrete state blocks with generous spacing (ANY at x=0.0, ERROR at x=7.6, targets at x=14.2)
    \node[draw, thick, rounded corners, fill=white, draw=gray!70!black, align=center, text width=2.8cm, minimum height=1.6cm, font=\small]
      (ANY) at (0.0, 0.6) {Any State $X \neq \text{ERROR}$\\[2pt]\footnotesize STANDBY, COLLECTING,\\\footnotesize READY, RUNNING, PAUSED};
    \node[state, fill=stErr] (E) at (7.6, 0.6) {ERROR\\[-1pt]\footnotesize\mdseries R:2Hz B:OFF};
    \node[state, fill=stPause, minimum size=2.1cm] (P) at (14.2, 2.2) {PAUSED};
    \node[draw, thick, rounded corners, fill=stIdle, align=center, text width=3.4cm, minimum height=1.5cm, font=\small\bfseries]
      (I) at (14.2, -1.2) {\textbf{STANDBY /}\\[2pt]\textbf{COLLECTING /}\\[2pt]\textbf{READY}};

    \draw[tr] (ANY) -- node[lbl, above=3pt, pos=0.45]{T17: FAULT} (E);
    \draw[tr] (E) to[out=65, in=115, looseness=4.0] node[lbl, above=2pt]{T18/T19: 1\,s Tick [Cycle]} (E);
    \draw[tr] (E.north east) -- node[lbl, sloped, above=3pt, pos=0.45]{T20: CLEARED [$t > 0$]} (P.west);
    \draw[tr] (E.south east) -- node[lbl, sloped, below=3pt, pos=0.45]{T21: CLEARED [$t = 0$]} (I.west);
    \node[font=\footnotesize, align=center, text width=13.0cm, fill=white, draw=gray!50, rounded corners, inner sep=4pt] at (7.1, -3.2)
      {Safety Lock-Out: RUN, PAUSE, STOP, and Coins are rejected during ERROR mode.\\
       Active cycle condition: Saved pre-fault state $X \in \{\text{RUNNING}, \text{PAUSED}\}$.};
  \end{tikzpicture}%
  }
  \caption{Formal state transition diagram: Asynchronous safety fault handling and context restoration}
  \label{fig:fsm-error}
\end{figure}

\subsection{State Transition Table}

\begingroup
\small
\setlength{\tabcolsep}{3pt}
\begin{longtable}{L{0.8cm}L{2.7cm}L{1.9cm}L{2.5cm}L{2.3cm}L{4.1cm}}
  \caption{Complete and deterministic state transition table}\label{tab:transitions}\\
  \toprule
  \textbf{ID} & \textbf{Current State} & \textbf{Input Event} & \textbf{Guard Condition} & \textbf{Next State} & \textbf{Action / Output Side Effects} \\
  \midrule
  \endfirsthead
  \toprule
  \textbf{ID} & \textbf{Current State} & \textbf{Input Event} & \textbf{Guard Condition} & \textbf{Next State} & \textbf{Action / Output Side Effects} \\
  \midrule
  \endhead
  \bottomrule
  \endfoot
  T01 & \st{STANDBY}    & COIN($v$) & $b+v < 50$ & \st{COLLECTING} & $b \mathrel{+}= v$; RLED \code{ON}, BLED \code{OFF} \\
  T02 & \st{STANDBY}    & COIN($v$) & $b+v \ge 50$ & \st{READY}      & $b \mathrel{+}= v$; RLED \code{OFF}, BLED \code{ON} \\
  T03 & \st{COLLECTING} & COIN($v$) & $b+v < 50$ & \st{COLLECTING} & $b \mathrel{+}= v$; $n := 0$; $w := 0$ \\
  T04 & \st{COLLECTING} & COIN($v$) & $b+v \ge 50$ & \st{READY}      & $b \mathrel{+}= v$; RLED \code{OFF}, BLED \code{ON} \\
  T05 & \st{COLLECTING}, \st{READY} & STOP & $n = 0$ & (same) & $n := 1$; $w := 1500\,\text{ms}$ \\
  T06 & \st{COLLECTING}, \st{READY} & STOP & $n = 1 \land w > 0$ & \st{STANDBY} & Refund deposit $b$; $b := 0$; $n := 0$; $w := 0$ \\
  T07 & \st{READY}      & COIN($v$) & -- & \st{READY}      & $b \mathrel{+}= v$ (surplus deposit forfeited on RUN) \\
  T08 & \st{READY}      & RUN       & -- & \st{RUNNING}    & $b := 0$ (no refund); $t := T$; door \code{LOCKED}; motor \code{AGITATE} \\
  T09 & \st{RUNNING}    & PAUSE     & -- & \st{PAUSED}     & motor, pump \code{OFF}; BLED \code{ON}; door stays \code{LOCKED}; \textbf{$t$ continues} \\
  T10 & \st{RUNNING}    & TICK$_{1\text{s}}$ & $t > 1$ & \st{RUNNING} & $t \mathrel{-}= 1$; motor \code{AGITATE} if $t > 300$, else \code{SPIN} + pump \\
  T11 & \st{RUNNING}    & TICK$_{1\text{s}}$ & $t = 1$ & \st{STANDBY} & $t := 0$; all actuators \code{OFF}; door \code{UNLOCKED}; RLED \code{ON} \\
  T12 & \st{RUNNING}, \st{PAUSED} & STOP & $n = 0$ & (same) & $n := 1$; $w := 1500\,\text{ms}$; washing/timer unaffected \\
  T13 & \st{RUNNING}, \st{PAUSED} & STOP & $n = 1 \land w > 0$ & \st{STANDBY} & Force abort: all \code{OFF}; unlock door; RLED \code{ON}, BLED \code{OFF} \\
  T14 & \st{PAUSED}     & RUN       & -- & \st{RUNNING}    & Resumes: BLED $1\,\text{Hz}$; restore motor by phase; $t$ \emph{preserved} \\
  T15 & \st{PAUSED}     & TICK$_{1\text{s}}$ & $t > 1$ & \st{PAUSED}     & $t \mathrel{-}= 1$ (monotonic countdown during pause) \\
  T16 & \st{PAUSED}     & TICK$_{1\text{s}}$ & $t = 1$ & \st{STANDBY} & $t := 0$; all actuators \code{OFF}; door \code{UNLOCKED}; RLED \code{ON} \\
  T17 & Any $\neq$ \st{ERROR} & FAULT & -- & \st{ERROR} & Save state; motor, pump \code{OFF}; door \code{UNLOCKED}; RLED $2.0\,\text{Hz}$ \\
  T18 & \st{ERROR}      & TICK$_{1\text{s}}$ & cycle $\land t > 1$ & \st{ERROR} & $t \mathrel{-}= 1$ (wall-clock budget expires during fault) \\
  T19 & \st{ERROR}      & TICK$_{1\text{s}}$ & cycle $\land t = 1$ & \st{ERROR} & $t := 0$; saved state $\gets$ \st{STANDBY} \\
  T20 & \st{ERROR}      & CLEARED   & cycle $\land t > 0$ & \st{PAUSED} & $f := 0$; BLED \code{ON}; user presses RUN to resume \\
  T21 & \st{ERROR}      & CLEARED   & otherwise & restored & $f := 0$; state chosen from $b$: \st{STANDBY}, \st{COLLECTING}, or \st{READY} \\
  T22 & Any             & TICK$_{1\text{ms}}$ & $w > 0$ & (same) & $w \mathrel{-}= 1$; if $w = 0$ then $n := 0$ (single press discarded) \\
\end{longtable}
\endgroup

\subsection{Determinism and Completeness Theorem}
\label{sec:determinism-proof}

\paragraph{Theorem 4.1 (Completeness and Determinism of EFSM $\mathcal{M}$):}
\textit{The state transition function $\delta: (\mathcal{S} \times \mathcal{C}) \times \Sigma \to (\mathcal{S} \times \mathcal{C})$
is mathematically complete and strictly deterministic:
\begin{equation}
  \forall (s, c) \in \mathcal{S} \times \mathcal{C}, \quad \forall e \in \Sigma, \quad \big| \delta((s, c), e) \big| = 1
\end{equation}
That is, for every possible state-context pair $(s, c)$ and every arriving event $e$, exactly one
valid next state and context mutation exists. Undefined behavior, nondeterministic branching, and
deadlocks are mathematically impossible.}

\paragraph{Proof of Determinism and Deadlock Freedom:}
To prove determinism and exhaustiveness, the entire Cartesian product space of qualitative states
and input events $|\mathcal{S}| \times |\Sigma| = 6 \times 10 = 60$ combinations is evaluated
against the implementation in Table~\ref{tab:event-matrix}:
\begin{enumerate}
  \item \textbf{Handled Transitions (22 Rules):} Transitions $\text{T01}$ through $\text{T22}$ partition
        the active operating domain. Every transition with guard conditions (e.g., $b+v < 50$ vs $b+v \ge 50$,
        or $t > 1$ vs $t = 1$, or $n = 0$ vs $n = 1 \land w > 0$) is formed by mutually exclusive and
        exhaustive predicates ($P \lor \neg P \equiv \text{True}$ and $P \land \neg P \equiv \text{False}$).
        Hence, no two transitions can fire simultaneously from the same configuration.
  \item \textbf{Filtered Non-Applicable Events (38 Combinations):} Any $(s, e)$ pair not explicitly defined
        in Table~\ref{tab:transitions} (e.g., \code{RUN} or \code{PAUSE} in \st{STANDBY}, or coin insertion
        during active washing \st{RUNNING}) is governed by the universal default filter rule:
        \begin{equation}
          \delta((s, c), e) = (s, c) \quad \text{and} \quad \omega_{\text{Mealy}}((s, c), e) = \emptyset
        \end{equation}
        The transition dispatcher cleanly filters the event, discarding the input
        with zero mutation of state or actuator outputs.
  \item \textbf{Deadlock Freedom:} The state graph contains no sink states. From any state $s \in \mathcal{S}$,
        either an internal timeout ($\text{TICK}_{1\text{s}}$), an explicit user action ($\text{STOP} \times 2$),
        or a fault clearance transition ($\text{CLEARED}$) guarantees a directed transition path back to
        the initial state $s_0 = \text{\st{STANDBY}}$.
\end{enumerate}
Therefore, $\mathcal{M}$ is complete, deadlock-free, and deterministic. \hfill $\blacksquare$

\begin{table}[H]
  \centering
  \caption{Exhaustive Event-Acceptance Matrix ($6 \times 10 = 60$ Combinations)}
  \label{tab:event-matrix}
  \small
  \resizebox{\linewidth}{!}{%
  \begin{tabular}{lcccccccccc}
    \toprule
    \textbf{State} & $\text{COIN}_{10}$ & $\text{COIN}_{20}$ & $\text{COIN}_{50}$ & \code{RUN} & \code{PAUSE} & \code{STOP} & $\text{TICK}_{1\text{s}}$ & $\text{TICK}_{1\text{ms}}$ & \code{FAULT} & \code{CLEARED} \\
    \midrule
    \st{STANDBY}    & T01/02 & T01/02 & T01/02 & \text{ign} & \text{ign} & \text{ign} & \text{ign} & T22 & T17 & \text{ign} \\
    \st{COLLECTING} & T03/04 & T03/04 & T03/04 & \text{ign} & \text{ign} & T05/06     & \text{ign} & T22 & T17 & \text{ign} \\
    \st{READY}      & T07    & T07    & T07    & T08        & \text{ign} & T05/06     & \text{ign} & T22 & T17 & \text{ign} \\
    \st{RUNNING}    & \text{ign} & \text{ign} & \text{ign} & \text{ign} & T09        & T12/13     & T10/11 & T22 & T17 & \text{ign} \\
    \st{PAUSED}     & \text{ign} & \text{ign} & \text{ign} & T14        & \text{ign} & T12/13     & T15/16 & T22 & T17 & \text{ign} \\
    \st{ERROR}      & \text{ign} & \text{ign} & \text{ign} & \text{ign} & \text{ign} & \text{ign} & T18/19 & T22 & \text{ign} & T20/21 \\
    \bottomrule
  \end{tabular}}
  \vspace{3pt}
  \raggedright\scriptsize\emph{Note:} \text{ign} denotes deterministic event filtering without state or context mutation.
\end{table}

\subsection{Moore Output Function Mapping}

\begin{table}[H]
  \centering
  \caption{Static Moore Output Function Mapping $\omega_{\text{Moore}}(s)$}
  \label{tab:moore}
  \small
  \resizebox{\linewidth}{!}{%
  \begin{tabular}{lcccccc}
    \toprule
    \textbf{State} & \textbf{RLED} & \textbf{BLED} & \textbf{Drum Motor} & \textbf{Drain Pump} & \textbf{Door Deadbolt} & \textbf{Countdown Timer} \\
    \midrule
    \st{STANDBY}    & Solid \code{ON}  & \code{OFF}       & \code{OFF}            & \code{OFF} & \code{UNLOCKED} & Inactive ($t=0$) \\
    \st{COLLECTING} & Solid \code{ON}  & \code{OFF}       & \code{OFF}            & \code{OFF} & \code{UNLOCKED} & Inactive ($t=0$) \\
    \st{READY}      & \code{OFF}       & Solid \code{ON}  & \code{OFF}            & \code{OFF} & \code{UNLOCKED} & Inactive ($t=0$) \\
    \st{RUNNING}    & \code{OFF}       & Blinking $1\,\text{Hz}$ & \code{AGITATE}/\code{SPIN} & \code{ON} in \code{SPIN} & \code{LOCKED}   & Monotonic decrement \\
    \st{PAUSED}     & \code{OFF}       & Solid \code{ON}  & \code{OFF}            & \code{OFF} & \code{LOCKED}   & \textbf{Monotonic decrement} \\
    \st{ERROR}      & Blinking $2\,\text{Hz}$ & \code{OFF}       & \code{OFF}            & \code{OFF} & \code{UNLOCKED} & Active if cycle interrupted \\
    \bottomrule
  \end{tabular}}
\end{table}

